\scanpage{9}
\mspage{4}
to which we shall restrict ourselves in what follows -- we thus have canonical homomorphisms
\[
(6)\qquad B_K\longrightarrow B_{\mathbb Q}\qquad\Pi_K\longrightarrow\Pi_{\mathbb Q}
\]
which are made explicit, when we have chosen an object of $\Pi_K$, that is, a $\bar K/K$, hence a $\overline{\mathbb Q}/\mathbb Q$, by a homomorphism of profinite groups
\[
(7)\qquad E_{\bar K/K}\longrightarrow\Gamma_{\overline{\mathbb Q}/\mathbb Q}.
\]
In what follows, we shall always regard $B_K$, $\Pi_K$ or $E_{\bar K/K}$ as endowed with this additional structure -- \unsure{it is} the ``\uline{arithmetic}'' \deleted{\unsure{f}}morphisms (of topoi, of groupoids, or of profinite groups), will dominate the situation. !

Particular interest attaches to the kernel of (7), which I denote by $\pi_{\bar K/K}$ -- it is called the ``geometric part'' of the Galois group $E_{\bar K/K}$, as opposed to the quotient
\[
E_{\bar K/K}/\pi_{\bar K/K}=\Gamma^{\deleted{\unsure{a}}}_{\bar K/K}\hookrightarrow\Gamma_{\overline{\mathbb Q}/\mathbb Q},
\]
\margin{we shall denote by $\Gamma=\Gamma^{\deleted{\unsure{a}}}_{\bar K/K}$ this ``arithmetic part''}
which I call its
\margin{8}
